\section{Numerical illustrations and further questions}
\label{sec:numerical}

The illustrations in this section were recomputed directly from the
reference-sphere formula \eqref{eq:reference_energy}, using ordinary
floating-point arithmetic.  They make the geometry of the proposed
classification visible; they are not inputs to Assertions C1--C6.
The accompanying source bundle contains the plotting code, sampled
values, quadrature orders, and the rational slice records used in
Figure~\ref{fig:branches}.

\subsection{Branches and meridian sections}

On the axis, rotational symmetry removes the azimuth integration
exactly.  We located the positive root at 69 parameters using a
768-node Gauss--Legendre rule, and compared it with a 384-node rule.
The largest change in the computed root was $9.87\times10^{-7}$.
The prolate branch was computed at 65 parameters with a
$192\times384$ polar--azimuth rule.  All roots were obtained by
bracketing and a scalar root solver, away from the transition
parameters.  The discrepancy between two quadrature resolutions
is a diagnostic, not a rigorous error estimate.

\begin{figure}[tbp]
\centering
\includegraphics[width=\linewidth]{figures/branches.pdf}
\caption{Numerically computed noncentral branches in reference
coordinates.  Left: the two axial saddles on the oblate side.
Right: the radius of the single prolate stationary circle.
The seven marked vertical intervals on the right are the exact
rational slice enclosures recorded in the window package; their
widths are mostly below the plotting resolution.  They retain the
evidence status of that package.  Open endpoint markers and dotted
connections indicate the classified limits, rather than additional
numerical root evaluations.  The centre persists throughout.}
\label{fig:branches}
\end{figure}

The energy sections in Figure~\ref{fig:energy} use a
$128\times256$ sphere rule and a $99\times99$ meridian grid,
retaining the 7,525 points with $q_1^2+q_3^2\leq0.98^2$.
The small omitted outer rim prevents a finite-resolution image from
suggesting that the boundary trace has been established by sampling.
Each panel has its own energy scale.  In the prolate circle regime,
two crosses in a meridian plane are intersections of one connected
stationary circle, not two isolated stationary points.

\begin{figure}[tbp]
\centering
\includegraphics[width=.96\linewidth]{figures/energy_sections.pdf}
\caption{Reference meridian sections of $\widehat E_\lambda$ for
$\lambda=0.3,0.52,3,6$.  A white disc marks the centre; white crosses
mark the noncentral stationary intersections.  The four panels
show, respectively, a central saddle, a central minimum with an
axial saddle pair, a central minimum with a saddle circle, and a
central saddle of index two.  The circle drawn around each panel
is the reference boundary; the energy is sampled only up to
radius $0.98$.  These are non-interval numerical illustrations.}
\label{fig:energy}
\end{figure}

\subsection{The size of the recorded sign margins}

The constants $16189/2^{36}$ and $3341881/2^{36}$ are small
positive residual margins from interval enclosures over entire
parameter and spatial boxes.  They are not estimates of the
actual global minimum of $V$, and the reporting unit $2^{-36}$
does not measure the full enclosure error.  Dependence between
repeated variables, angular range overestimation, and outer-box
width can all consume a large part of a sign margin
\citep{Rump2010}.  For the two vertical windows of C4, the reported whole-box margins for $E_{zz}$ and $E_z$ are about $1.75\times10^{-4}$ and $4.86\times10^{-5}$. A separate numerical review reported sampled minima near $1.07$ and $0.32$ in the respective windows. The several-order gap illustrates enclosure loss; the sampled minima supply no lower bound on the continuum. A useful scale comparison is provided by direct
point values at $\rho=z=0.2$:
\begin{center}
\begin{tabular}{rrr}
\toprule
$\lambda$&$V(0.2,0.2;\lambda)$&$H(0.2,0.2;\lambda)$\\
\midrule
1&1.261673&1.261673\\
1.5&1.810885&0.883353\\
2&2.053848&0.593574\\
3&2.169992&0.242475\\
5&2.025279&$-0.042247$\\
\bottomrule
\end{tabular}
\end{center}
Doubling the rule from $128\times256$ to $256\times512$
changes these point values by at most $1.65\times10^{-14}$.
This agreement does not certify the digits and does not supply
a continuum lower bound.  It only illustrates why a very small
recorded lower endpoint need not indicate a nearly vanishing
derivative at a typical interior point.

\subsection{A possible extension of axial concavity}

The analytic concavity lemma proves a uniform estimate for
$0<\lambda\leq1/32$.  The following stronger sign statement
would simplify much of the remaining oblate analysis.

\begin{conjecture}\label{conj:concavity}
For every $0<\lambda<\cO$ and $|q|<1$,
\[
 \partial_{zz}\widehat E_\lambda(q)<0.
\]
\end{conjecture}

A numerical check with a $96\times192$ sphere rule sampled
473 points in the nonnegative meridian quarter-disc of radius
$0.96$ for each of
$\lambda=1/16,1/8,1/4,0.35,0.4,0.4079$.
All sampled values were negative.  The maxima in these finite
samples were approximately
$-0.223524,-0.367927,-0.328099,-0.138342,-0.019734$,
and $-0.000147$, respectively.
A further chord-Hessian quadrature at 17 evenly spaced meridian angles on each of eight near-boundary arcs gave strictly negative sample values throughout. For $\lambda=1/128,1/64,1/32$ we used radius $0.9999$. For $\lambda=0.25,0.35,0.4,0.4079,0.42$ we used $0.9995$. The least negative sampled values on the two sets were about $-0.031$ and $-0.487$. These are point diagnostics. In particular, $0.42>\cO$ lies outside the conjecture and the centre there already has positive axial curvature.

The assertion that the maximum always occurs at the centre would
be stronger and is not suggested by all the samples: at
$\lambda=1/8$ the value at $(\rho,z)=(0.88,0)$ is about
$-0.367927$, larger than the central value $-0.383316$.
For the four sampled parameters from $1/4$ through $0.4079$,
the largest sampled value does occur at the centre.
Conjecture~\ref{conj:concavity} does not include a claim about
the location of the maximum, and no computation in this section
is used to eliminate an existing proof obligation.

\subsection{A concrete triaxial question}

Fix $\lambda$ strictly inside $(\bP,\cP)$ and break rotational
symmetry by replacing $M_\lambda$ with
\[
 M_{\lambda,\epsilon}
 =\diag(e^\epsilon,e^{-\epsilon},1)M_\lambda.
\]
Its determinant remains one.  Conditional on the stationary-circle
classification, the normal Hessian along the unperturbed circle
is invertible.  The implicit function theorem therefore solves
the two normal critical equations near that circle and reduces
the remaining equation to a periodic angular energy
$\Psi_\epsilon(\phi)$.
The first variation is linear in
$T=\diag(1,-1,0)$.  Rotation of this tensor through $\phi$
produces its cosine and sine components of angular weight two;
reflection across $q_2=0$ removes the sine component at
$(\rho_*,0,0)$.  The first-order contribution of the normal
correction vanishes because the unperturbed gradient is zero.
Consequently, with
$A(\lambda)=\left.\partial_\epsilon
\widehat E_{\lambda,\epsilon}(\rho_*,0,0)\right|_{\epsilon=0}$,
the reduced expansion is
\[
 \Psi_\epsilon(\phi)
 =\Psi_0+\epsilon A(\lambda)\cos(2\phi)
   +O_{C^2(S^1)}(\epsilon^2),
\]
provided $A(\lambda)\ne0$.
For sufficiently small nonzero $\epsilon$, the four simple
zeros of the leading angular derivative persist and exhaust
the nearby critical points.  Coordinate reflections, inherited
by the unique normal solution, place them exactly at
$\phi=k\pi/2$, in two opposite principal-axis pairs.
Their angular curvatures alternate in sign.  A Schur complement
of the normal Hessian gives a congruent block decomposition
with the reduced angular curvature.  Since the two normal
curvatures keep opposite signs, two of the points have Morse
index one and two have index two.

There is also an exact angular equation. Write the energy as $F(C,q)$ for the positive definite matrix $C=L^TL$, and define the symmetric matrix $G$ by $D_CF[U]=\operatorname{tr}(GU)$. Put $T_{ij}=2\ell_i\ell_jG_{ij}$ in a principal-axis frame with semiaxes $\ell_i$. For every rotation $R$, $F(RCR^T,Rq)=F(C,q)$ by change of the reference integration variable. Differentiating a rotation in the $(1,2)$ plane yields
\[
 q_1E_2-q_2E_1=(\ell_2/\ell_1-\ell_1/\ell_2)T_{12}.
\]
The differential of the displayed diagonal deformation similarly gives $A=T_{11}-T_{22}$ at the unperturbed circle point on the first axis. The exact identity also applies at nonzero deformation.

For fixed shape near $\cP$, the shape derivative on the first axis is smooth and even in $\rho$. Its constant term vanishes by interchange of the two equatorial axes. Hence $A=O(\rho_*^2)$ as the circle shrinks. The natural quantity to estimate uniformly is $A/\rho_*^2$. Ordinary quadrature gives $0.6744,0.5744,0.4879,0.4399$ at $\lambda=2.1,3,4,4.7$ respectively. A positive uniform lower bound on the full branch remains an open estimate here. The local triaxial study cited in the introduction concerns the neighbourhood of the prolate centre transition. Neither these four samples nor the exact angular equation by itself excludes other critical points over the full triaxial family.

\section{Availability, reproducibility, and status of the draft}
\label{sec:availability}

The computational references identify fixed commits and paths in
the public repositories \texttt{cotaxxxx/bg-spheroid-family} and
\texttt{cotaxxxx/bg-oblate-spheroid}.
The earlier background material is in the public repository
\texttt{cotaxxxx/basepoint-geometry}.  The analytic flat-end
argument and the axial part of the prolate-tail argument are
expanded here from their fixed research records
\citep{data-oblate-analytic,data-prolate-axis}.
The sphere proof is given in Section~\ref{sec:sphere} and does
not require an unpublished spherical manuscript.
Record V24 (16 November 2026): \texttt{cotaxxxx/bg-spheroid-family},
commit \texttt{127be52046e9498d1306fe3d35c3e02e91c636de},
file \href{https://github.com/cotaxxxx/bg-spheroid-family/blob/127be52046e9498d1306fe3d35c3e02e91c636de/runs/2026-11-16/D_OB_GLOBAL_CLOSED_V24/JUDGE_RULING_D_OB_GLOBAL_CLOSED_20261116.md}{\path{JUDGE_RULING_D_OB_GLOBAL_CLOSED_20261116.md}}
(Japanese working record). An exact copy of this primary record is
included in the accompanying source bundle under \path{provenance/V24/}.
The accompanying project-level AI audit checked the assembly argument
and source references; it did not independently rerun the computations.

For a claim-by-claim reproduction map, see
\path{CURRENT_REPRODUCTION_v0_5_4.md} in the source bundle. It lists
C1--C6 with pinned input packages, producer and checking commands,
ledger paths, and the extent of independent reevaluation, separately
from the retained historical \path{REPRODUCTION.md}. The new map does
not claim that every listed original producer has been rerun or that
its original analytic kernel has been independently rederived.

The accompanying revision bundle contains the manuscript sources,
bibliography, vector figures, figure data and reproduction scripts,
and a map of the fixed computational sources.  It is not the
complete certified archive of every upstream computation.  No
archival DOI for such a consolidated proof package is claimed
in this revision.  A fixed release should include the exact
input boxes, all selected final leaves, coefficient tables,
producer and checker sources, build environment, execution
commands, and file hashes, with a direct map from C1--C6 to
the corresponding accepted records.

There are three distinct obligations: correctness of the analytic
kernel, correct outward arithmetic and error bounds in its
implementation, and completeness of the domain cover used by
the classification.  A file hash establishes identity, not any
of these mathematical properties.  Reproduction with the original
kernel and reevaluation with a separately derived kernel address
different failure modes.  Full independent numerical replay is
not asserted here; its current scope is stated in
Section~\ref{subsec:computational-workload}.
The supplement gives the replay commands, fixed locations and selection rules. The source bundle contains the independent interval calculations and their integer coverage checks.

\paragraph{Use of AI tools.}
ChatGPT (OpenAI) assisted with exploratory derivations, numerical
code and verification scripts, English-language wording proposals,
revision and typesetting. Claude (Anthropic) assisted with separate
project-level AI audits of argument assembly, source references and
manuscript differences. These AI audits are not an independently
implemented replay of all computational inputs or external human peer
review. The provenance of first-draft English paragraphs as opposed
to subsequent translation or copyediting has not yet been fully
reconstructed; an account restricted to ``copyediting'' is therefore
not asserted. The source bundle includes an AI-provenance review log.
No compliance with the first-draft restriction of a prospective
journal is claimed. Author verification of every suggested claim,
reference, code contribution and the final disclosure is pending.
